\begin{textsample}{POS Dim 7 – global\_north\_2024\_09 – Score 2.00 – global\_north\_2024\_09/114400858.txt}  \label{ex:f7_pos_435}
Canada abstains from UN motion calling on Israel to end occupation of Gaza, West Bank Bob Rae, ambassador of Canada to the United Nations, speaks while holding a copy of the United Nations charter at the general assembly hall, Wednesday, Feb. 23, 2022, at United Nations Headquarters.
Canada abstained today in a high-profile United Nations vote demanding that Israel end its"unlawful presence"in Gaza and the occupied West Bank within a year.
THE CANADIAN PRESS/AP/John Minchillo By Dylan Robertson, The Canadian Press Posted September 18, 2024 1:33 pm.
Last Updated September 18, 2024 7:10 pm.
OTTAWA -- Canada abstained Wednesday from a high-profile United Nations vote demanding that Israel end its"unlawful presence"in the Gaza Strip and the occupied West Bank within a year.
The decision follows a shift in how Ottawa votes at the UN, and it alarmed a prominent Israel advocacy group while earning criticism from two Liberal MPs."We can not support a \textbf{resolution} where one party, the State of Israel, is @ @ @ @ @ @ @ @ @ @ ambassador to the United Nations, Bob Rae, told the General Assembly.
The State of Palestine brought the non-binding motion, which passed 124-14 ; Canada was among 43 \textbf{abstentions}.
The \textbf{resolution} was based on a non-binding ruling by the top United Nations court in July that said Israel ' s presence in the Palestinian territories is unlawful and must end.
The motion pertained to Israel ' s decades of occupation in the West Bank, as well as its war against Hamas in Gaza.
It comes as the first anniversary of the war approaches and as violence in the West Bank reaches new highs.
Riyad Mansour, the Palestinian UN ambassador, called the vote a turning point"in our struggle for freedom and justice.""It sends a clear message that Israel ' s occupation must end as soon as possible and that the Palestinian people ' s right to self-determination must be realized,"he said."Instead of marking the anniversary of the Oct. 7 massacre by condemning Hamas and calling for the release @ @ @ @ @ @ @ @ @ @ Assembly continues to dance to the music of the Palestinian Authority,"Danon said.
Rae said Canada abstained despite agreeing that Israel is illegally occupying Palestinian territories, because the \textbf{resolution} calls for isolating Israel and contains language around boycotting the country, which he said will not lead to peace.
He said the violence in Israel and the territories it occupies is"a conflict where everyone in this room knows that many other states and non-state actors are also directly involved."While the \textbf{resolution} is not legally binding, the extent of its support reflects world opinion.
There are no \textbf{vetoes} in the General Assembly, unlike in the 15-member Security Council.
The \textbf{resolution} also demands the withdrawal of all Israeli forces and the evacuation of settlers from the occupied Palestinian territories"without delay."And it urges countries to impose sanctions on those responsible for maintaining Israel ' s presence in the territories and halt arms exports to Israel if they ' re suspected of being used there.
In addition, the \textbf{resolution} calls @ @ @ @ @ @ @ @ @ @ caused by its occupation and urges countries to take steps to prevent trade or investments that maintain Israel ' s presence in the territories.
Liberal MPs Anthony Housefather and Marco Mendicino said Canada should have voted against the"blatantly one-sided"\textbf{resolution}, arguing the \textbf{abstention} doesn't advance conditions for peace or adequately recognize Israel ' s right to exist.
The Centre for Israel and Jewish Affairs argued that Canada on Wednesday violated a statement Rae ' s delegation issued last November, which said Canada would vote down motions unfairly targeting Israel."Canada reiterates the importance of a fair-minded approach at the United Nations and will continue to vote ' no ' on \textbf{resolutions} that do not address the complexities of the issues or seek to address the actions and responsibilities of all parties, including the destructive role of terrorist organizations like Hamas, Palestinian Islamic Jihad and Hezbollah,"the Nov. 9, 2023 statement reads.
That statement was issued as Canada abstained on a motion calling for an"immediate, durable and sustained humanitarian truce"@ @ @ @ @ @ @ @ @ @ UN motions focusing on Israel, as the U.S. continues to do.
A month later, Canada broke with tradition by voting in favour of a UN motion seeking an"immediate ceasefire"without specifically condemning Hamas.
Reacting to Wednesday ' s vote, CIJA argued Canada is rewarding Hamas by abstaining and not voting against the motion."Canada broke its promise and put political expediency before principle,"the group wrote in a statement.
The International Court of Justice ruling that spawned the motion was a sweeping condemnation of Israel ' s rule over the lands it captured during the 1967 war.
It said Israel had no right to sovereignty over the territories and was violating international laws against acquiring the lands by force.
The motion asks UN Secretary-General Antonio Guterres to submit a report to the General Assembly within three months of putting the \textbf{resolution} in place,"including any actions taken by Israel, other states and international organizations, including the United Nations.""We fully abide by the decisions of @ @ @ @ @ @ @ @ @ @."I will implement any decision of the General Assembly in that regard."Meanwhile, Canada laid sanctions Wednesday against people with"roles in Hamas ' s financial network"as well as"extremist"Israeli settlers.

% matched lemmas: abstention, resolution, veto
\end{textsample}
